\documentclass[11pt,a4paper]{article}
\usepackage[T1]{fontenc}
\usepackage[utf8]{inputenc}
\usepackage{lmodern}
\usepackage[margin=27mm,headheight=15pt]{geometry}
\usepackage{amsmath,amssymb,amsthm,mathtools}
\usepackage{microtype}
\usepackage{booktabs,longtable,array}
\usepackage[dvipsnames]{xcolor}
\usepackage{xurl}
\usepackage{fancyhdr}
\usepackage{hyperref}
\hypersetup{colorlinks=true,linkcolor=MidnightBlue,urlcolor=MidnightBlue,
  citecolor=MidnightBlue,pdftitle={A computer-assisted construction of a c=8 Haagerup CFT},
  pdfauthor={Research manuscript},pdfsubject={Proof manuscript with explicit review obligations}}
\urlstyle{tt}
\setlength{\emergencystretch}{2em}
\setlength{\parskip}{0.35em}
\pagestyle{fancy}
\fancyhf{}
\fancyhead[L]{\small Central-charge-eight Haagerup construction}
\fancyhead[R]{\small Review draft}
\fancyfoot[C]{\thepage}
\newtheorem{theorem}{Theorem}[section]
\newtheorem{lemma}[theorem]{Lemma}
\newtheorem{proposition}[theorem]{Proposition}
\newtheorem{corollary}[theorem]{Corollary}
\newtheorem{claim}[theorem]{Claim}
\theoremstyle{definition}
\newtheorem{definition}[theorem]{Definition}
\newtheorem{remark}[theorem]{Remark}
\newtheorem{certificate}[theorem]{Certificate input}
\DeclareMathOperator{\Rep}{Rep}
\DeclareMathOperator{\Com}{Com}
\DeclareMathOperator{\FPdim}{FPdim}
\DeclareMathOperator{\Hom}{Hom}
\DeclareMathOperator{\End}{End}
\DeclareMathOperator{\Tr}{Tr}
\DeclareMathOperator{\wt}{wt}
\DeclareMathOperator{\Span}{span}
\DeclareMathOperator{\qdim}{qdim}
\DeclareMathOperator{\ch}{ch}
\DeclareMathOperator{\Aut}{Aut}
\renewcommand{\Vec}{\mathrm{Vec}}
\newcommand{\HH}{\mathcal H}
\newcommand{\CC}{\mathbb C}
\newcommand{\QQ}{\mathbb Q}
\newcommand{\ZZ}{\mathbb Z}
\newcommand{\one}{\mathbf1}
\newcommand{\artifact}[1]{\textnormal{\nolinkurl{#1}}}
% Breakable review markers require only standard LaTeX and xcolor.
% Their textual labels remain meaningful in a monochrome printout.
\newenvironment{reviewnote}[1]{\par\addvspace{8pt}\begingroup
  \color{Sepia}\noindent\rule{\linewidth}{0.7pt}\par\nobreak
  \noindent\textbf{REVIEW --- #1}\par\nobreak\small\noindent\ignorespaces
}{\par\noindent\rule{\linewidth}{0.7pt}\par\endgroup\addvspace{8pt}}
\newenvironment{criticalnote}[1]{\par\addvspace{8pt}\begingroup
  \color{BrickRed}\noindent\rule{\linewidth}{1pt}\par\nobreak
  \noindent\textbf{HIGH PRIORITY REVIEW --- #1}\par\nobreak\small\noindent\ignorespaces
}{\par\noindent\rule{\linewidth}{1pt}\par\endgroup\addvspace{8pt}}
\title{A computer-assisted construction of a\\[3pt]
  \texorpdfstring{\(c=8\)}{c=8} Haagerup CFT\\[5pt]
  \large Proof manuscript with explicit review obligations}
\author{Research manuscript}
\date{23 September 2026}

\begin{document}
\maketitle
\begin{abstract}
We describe an explicit conformal vertex subalgebra of the lattice VOA
\(V_{E_6\oplus A_2}\), generated by an \(A_{1,1}\otimes A_{1,39}\)
current algebra, the ambient stress tensor, and one three-term lattice field.
Exact finite certificates, together with a recursive strong-generation
argument, give a proposed computer-assisted proof that its current space
has dimension six. We present the subsequent regularity, branching,
modular-character and categorical arguments leading to a proposed
identification with the Haagerup center. The full unitary-CFT conclusion
depends on the analytic claims highlighted in red. Additional delicate
local deductions and imported proof components are marked in amber.
This document records a proof for review; it is not a report of external
expert validation or proof-assistant verification.
\end{abstract}

\begin{criticalnote}{Status of the main claim}
The computational closure certificates are complete. Their completion
does not mechanically verify the new analytic arguments used later in
the manuscript. In particular, the net-sector comparison, module
reconstruction/fullness, and annulus compatibility require specialist
review. The complete-unitarity and unitary ribbon-equivalence conclusions
are stated conditionally on those claims below. Earlier research notes
that call the entire identification ``complete'' should be read with
this explicit qualification. Reviews by additional agents are internal
reviews, not independent human expert validation.
\end{criticalnote}

\clearpage
\tableofcontents
\clearpage

\section{Statement, conventions, and proof dependencies}
\label{sec:statement}

Let
\[
 B=V_{E_6\oplus A_2},\qquad H=V_{E_8},\qquad B\subset H
\]
be the standard index-three lattice extension. The notation \(A_{1,k}\)
means the simple affine VOA \(L_k(\mathfrak{sl}_2)\). Fix the embedding
\(\mathcal K=A_{1,1}\otimes A_{1,39}\subset B\) described in
Section~\ref{sec:construction}. Let \(u\in B_2\) be the field in
\eqref{eq:u}, and define the \emph{actual generated} vertex algebra
\begin{equation}\label{eq:U}
 U=\langle\mathcal K,\omega_B,u\rangle\subset B.
\end{equation}
All products are computed in this lattice VOA. Thus associativity is
inherited from the ambient algebra, rather than assumed for an abstract
presentation inferred from character coefficients.
Throughout the normal-space and mode formulas, \(K=\mathcal K_1\)
denotes only the six-dimensional current space, not the full affine VOA.

\begin{proposition}[Computational properness conclusion]\label{prop:mainproper}
Accept the exact lattice identities and provenance recorded in
Section~\ref{sec:ledger}. The algebra \(U\) is strongly generated by
actual fields of weights at most four. Its weight-one space is
\[
 U_1=\mathfrak{sl}_2\oplus\mathfrak{sl}_2,\qquad \dim U_1=6.
\]
Consequently \(U\subsetneq B\), since \(\dim B_1=86\).
\end{proposition}
The finite certificates and the all-weight implication are presented in
Sections~\ref{sec:finite}--\ref{sec:properness}. This is the most directly
checkable part of the construction.

\begin{claim}[Full identification proposed by the manuscript]\label{claim:main}
With the additional local deductions and analytic claims marked for
review in Sections~\ref{sec:regularity} onward, \(U\) is completely
unitary and strongly rational, and
\begin{align}
 c(U)&=8,\qquad \Rep(U)\simeq Z(\HH_3),\label{eq:maincat}\\
 \ch_U(\tau)&=q^{-1/3}\bigl(1+6q+120q^2+956q^3+6010q^4+\cdots\bigr).
 \label{eq:vacchar}
\end{align}
The equivalence is an equivalence of unitary ribbon categories.
The full twelve-component character vector is the \(\gamma=0\)
Evans--Gannon vector. There is a unitary diagonal full CFT with
\(c_L=c_R=8\) and partition function
\(\sum_i|\chi_i(\tau)|^2\). A choice of holomorphic extension
\(H'\cong V_{E_8}\) has all-module category \(\HH_3\), and its
algebra-bimodule category is \(\HH_3\boxtimes\HH_3^{\mathrm{op}}\).
\end{claim}

Here \(H'\) is obtained by transporting the canonical \(\HH_3\)
Lagrangian algebra. It need not coincide with the originally embedded
extension \(H/U\); the proposed argument identifies that original
all-module category only up to the alternatives \(\HH_2,\HH_3\).
The full-CFT defects in Claim~\ref{claim:main} preserve \(U\), and need
not commute with the enlarged \(E_8\) chiral algebra.

\subsection{How the qualifications are used}
Amber \textbf{REVIEW} notes identify deductions whose complete verification
is not supplied by the finite closure computations, or whose long local
proofs are imported from companion notes. Red \textbf{HIGH PRIORITY REVIEW}
notes identify claims about which confidence is lower, especially operator
domains and the passage from conformal nets to complete VOA unitarity.
A highlighted claim is not being declared false. It is deliberately made
visible as a premise of the downstream conclusion.

\begin{center}
\begin{tabular}{@{}p{0.27\linewidth}p{0.65\linewidth}@{}}
\toprule
Part of the argument & Nature of its evidence \\
\midrule
Explicit fields and finite closure & Exact arithmetic over
\(\QQ(\sqrt{13},\sqrt{-3})\), saved recipes and full residuals.\\
Properness in all weights & Recursive algebraic theorem plus finite
certificates; no character upper-bound extrapolation.\\
Rationality and branching & Positive-stress and category arguments,
finite-top certificates, and cited external theorems.\\
Haagerup modular data & Rank exclusions, indicator and character
arguments; some long derivations remain companion proof inputs.\\
Complete unitarity and unitary identification & New analytic comparison
claims, followed by cited unitary classification and extension theorems.\\
\bottomrule
\end{tabular}
\end{center}

\section{The explicit lattice construction}\label{sec:construction}
Use simple lattice vectors indexed \(0,\ldots,7\). The Gram matrix
has diagonal entries two and off-diagonal entries minus one on
\[
 (0,1),(1,2),(2,3),(3,4),(2,5),(6,7),
\]
with all remaining off-diagonal entries zero. Fix the cocycle
\begin{equation}\label{eq:cocycle}
 \epsilon(a,b)=(-1)^{\sum_i a_i b_i+
                      \sum_{i>j}a_i b_jG_{ij}}.
\end{equation}
The first current factor corresponds to the root
\(a_0=(-1,-2,-3,-2,-1,-2,0,0)\). The second is the diagonal of the
principal \(A_1\) in \(A_5\), at level 35, and the principal \(A_1\)
in the external \(A_2\), at level four.

The exact coefficient convention uses unnormalized simple-root
Heisenberg oscillators and a torus gauge in which the principal raising
current is the sum of the corresponding simple-root vertices. The
lowering current is \(-\sum_i c_i e^{-\alpha_i}\), with
\[
 (c_0,\ldots,c_7)=(5,8,9,8,5,0,2,2).
\]
The torus rescales simple-root vertices by
\( (1/\sqrt5,1/\sqrt8,1/3,1/\sqrt8,1/\sqrt5,1,1/\sqrt2,1/\sqrt2)\).
In this convention define
\begin{align}
 u={}&e^{(0,0,0,0,0,-1,1,1)}\notag\\
 &+\frac{\sqrt{13}+3\sqrt{-3}}{20}
                   e^{(0,0,0,1,1,-1,0,0)}\notag\\
 &+\frac{\sqrt{13}-3\sqrt{-3}}{20}
                   e^{(1,1,0,0,0,-1,0,0)}.\label{eq:u}
\end{align}
The compact real form and PCT operation must be transported through
this gauge. Ordinary coefficient conjugation in the displayed gauge
is not a replacement for the ambient PCT map.

\begin{certificate}[Actual initial spaces]\label{cert:initial}
The saved exact construction gives a PCT-stable 120-dimensional
space \(S_2\subset U_2\), containing \(u,\omega_B\), the affine
vacuum descendants, two neutral primitive fields, and eight charged
current multiplets. The current multiplet dimensions are
\[
 (1,5),(1,9),(1,9),(1,13),(2,2),(2,6),(2,8),(2,14).
\]
These contribute 96 vectors, in addition to the 22 affine vacuum
vectors and two neutral fields. The exact generated grade-three
space \(S_3\subset U_3\) has dimension 956. Its normal subspace
\(N_3=K_{-1}S_2+TS_2\) has dimension 634, and
\[
 S_3=N_3+\Span\{a_0b:a,b\in S_2\},\qquad \dim(S_3/N_3)=322.
\]
The spaces are current-stable. These are dimensions of constructed
finite spaces, not initial assertions that they exhaust \(U_2,U_3\).
\end{certificate}
The source records are
\artifact{hi_cft_20260923_c8_exact_fields.json},
\artifact{hi_cft_20260923_c8_exact_unitary_certificate.json}, and
\artifact{hi_cft_20260923_c8_minimal_generators.json}, together with
the exact grade-three checkpoint bound into the final ledger.

\section{Finite closure and the all-weight criterion}\label{sec:finite}
Write \(Y(a,z)=\sum_{n\in\ZZ}a_nz^{-n-1}\) and \(T=L(-1)\).
A singular product has mode index \(n\geq0\). The zero mode \(a_0\)
is a derivation and commutes with \(T\).

\begin{theorem}[Recursive strong-generation criterion]\label{thm:recursive}
Let \(V\) be a nonnegatively graded vertex algebra with \(V_0=\CC\one\).
Choose finite homogeneous subspaces \(S_d\subset V_d\),
\(1\leq d\leq h\), and let \(\mathcal N\) be their differential
normally ordered span, allowing every parenthesization. Put
\(K=S_1\) and \(A=S_2\). Assume
\begin{enumerate}
\item \(k_ns,a_ns\in\mathcal N\) for every \(k\in K\), \(a\in A\),
\(s\in S_d\), and \(n\geq0\);
\item for \(d\geq3\),
\[
 S_d\subset A_0S_{d-1}+\mathcal N_d(S_1+\cdots+S_{d-1}).
\]
\end{enumerate}
Then \(\mathcal N\) is a vertex subalgebra. In particular, the chosen
spaces strongly generate the vertex algebra that they generate.
\end{theorem}

\begin{proof}
The first assumption, the derivation property and \([T,a_0]=0\)
give \(a_0\mathcal N\subset\mathcal N\) for \(a\in A\).
We first recall the filtered Wick observation. If all generator pairs
of total weight less than \(w\) have singular products in \(\mathcal N\),
the same is true for normal expressions of total weight less than \(w\).
At total weight \(w\), a derivative or a product with at least two
positive-weight factors also reduces to smaller total weights. Indeed,
\begin{align*}
 [x_\lambda :yz:]
  ={}&:[x_\lambda y]z:+:y[x_\lambda z]:
       +\int_0^\lambda [[x_\lambda y]_\mu z]\,d\mu.
\end{align*}
Each omitted positive-weight factor lowers the first two pair weights.
In a nested term, \(x_jy\) has weight \(\wt(x)+\wt(y)-j-1\), so
its pairing with \(z\) has total weight at most \(w-1\).
Sesquilinearity handles derivatives and skew symmetry handles a
composite on the left. These are identities inside \(V\).

Now induct lexicographically on
\((\wt(s)+\wt(t),\wt(s))\) for generator pairs \(s,t\).
Left weight one or two is the first assumption. Otherwise the second
assumption reduces the left input to \(s=a_0b\), with
\(\wt(b)=\wt(s)-1\), plus normal remainders. The latter are covered
by the filtered Wick observation. Use
\begin{equation}\label{eq:recursive-jacobi}
 (a_0b)_nt=a_0(b_nt)-b_n(a_0t).
\end{equation}
The first term has smaller total pair weight before applying \(a_0\),
which preserves \(\mathcal N\). The known normal expression for
\(a_0t\) has weight \(\wt(t)+1\). Every derivative or composite term
reduces total pair weight in the second term. The only remaining
possibility is a single undifferentiated generator of weight
\(\wt(t)+1\); its pairing with \(b\) has the original total weight
and smaller left weight. It is therefore an earlier induction case.
If \(\wt(t)=h\), this single-generator possibility is absent.

All generator singular products close. Wick induction then gives closure
for all normal expressions. Translation and normal multiplication
preserve \(\mathcal N\) by definition, so all negative modes close too.
\end{proof}
For the Wick identities see
\href{https://arxiv.org/html/math-ph/0312042#S2}{De Sole--Kac, Section 2}.
The full local derivation is
\artifact{hi_cft_20260923_next_structure_recursive_strong_certificate.md}.

\begin{reviewnote}{F2: the all-weight implication is an algebraic proof}
The filtered Wick reduction and the lexicographic induction above are
new deductions connecting finitely many gates to closure in every weight.
Their correctness is not verified by checking more low-weight products.
In particular, a reviewer should check the reduction of composite normal
expressions and the single-generator case at unchanged total weight.
The finite certificates below supply the hypotheses; they do not replace
this argument.
\end{reviewnote}

\subsection{The finite spaces and prerequisite controls}
Choose actual quotient lifts \(P_3\subset S_3\) and an actual triplet
\(P_4\subset U_4\), with joint current highest weight \((0,2)\), and set
\begin{align}
 N_4&=K_{-1}S_3+TS_3+(S_2)_{-1}S_2,&W_4&=N_4+P_4,\label{eq:normal4}\\
 N_5&=K_{-1}W_4+TW_4+(S_2)_{-1}S_3.&&\label{eq:normal5}
\end{align}
These denote full differential normal spaces, using the lower closure
identities to reorder expressions. A computational solver may use a
proper subset of actual words in these spaces.

The triplet is obtained from an actual \(a_0p\) by current-zero-mode
projectors and lowering. It satisfies
\(P_4\subset (S_2)_0S_3+N_4\). At ambient grade four, the two
Cartan norms \(2,78\) bound highest weights by \(4,24\), respectively.
Consequently finite polynomials in \(FE\) give exact projectors.
For an input of weight \((0,2)\), one may use
\[
 \prod_{\ell=1}^{2}\left(1-\frac{F_LE_L}{\ell(\ell+1)}\right),
 \qquad
 \prod_{\ell=2}^{12}\left(1-\frac{F_RE_R}{\ell(\ell+1)-2}\right).
\]
The raising operator acts first in every \(FE\) factor.

\begin{certificate}[Lower identities and the actual triplet]\label{cert:lower}
The lower-closure artifact contains 3,970 completed exact identities.
The actual \(P_4\) vectors satisfy 18 current-zero-mode identities and
six independently checked inclusions
\[
 L(1)P_4\subset S_3,\qquad L(2)P_4\subset S_2.
\]
These direct Virasoro checks are available before the weight-four gate.
\end{certificate}
The two gates left for the all-weight argument are
\begin{align}
 (S_2)_0S_3&\subset W_4 &&\text{(Gate 4)},\label{eq:gate4}\\
 (S_2)_0P_4&\subset N_5 &&\text{(Gate 5)}.\label{eq:gate5}
\end{align}

\section{Gate 4: exact seeds, raising, and singlet recovery}\label{sec:gate4}
Put \(P=S_3/N_3\), of independently known dimension 322. For a fixed
joint-highest primitive field \(a\in S_2\), the map
\[
 f_a:P\longrightarrow B_4/W_4,\qquad [p]\longmapsto[a_0p]
\]
is well defined because the lower relations give \(a_0N_3\subset N_4\).
Current covariance implies that its kernel is stable under both commuting
raising operators \(E_L,E_R\).

\begin{lemma}[Raising saturation]\label{lem:raising}
If \(C_a\subset\ker f_a\) and
\(P=E_LP+E_RP+C_a\), then \(f_a=0\).
\end{lemma}
\begin{proof}
On \(M=P/\ker f_a\) the hypothesis says
\(M=E_LM+E_RM\). The commuting ideal \((E_L,E_R)\) acts nilpotently
because all current weights are bounded. Iterating gives
\(M=(E_L,E_R)^rM=0\) for sufficiently large \(r\).
\end{proof}

The exact raising coordinates comprise 644 images. Their rank is 294.
For each of nine independent highest left fields, the completed seed
sets and these images have rank 322. A nonzero minor at the good split
prime \(p=1{,}000{,}003\) proves a nonzero characteristic-zero minor.
This proves spanning because the exact ambient quotient dimension is
already 322. It does not give an upper rank for an unknown weight-four
normal space.

The original exact run supplies 2,145 seed identities. The saturation
calculation asks for 130 further inputs; 118 have direct exact residual
certificates. The remaining twelve outputs have weight \((0,0)\).
An exact replay of the known kernel and twelve additional identities
at output weight \((0,2)\) shows that both raising currents kill each
of these twelve remaining classes modulo \(W_4\).

\begin{lemma}[Recovery from a joint-singlet projection]\label{lem:singlet}
Let \(t\in B_4\) have weight \((0,0)\), with
\(E_Lt,E_Rt\in W_4\). Define
\begin{equation}\label{eq:projector}
 \Pi=\prod_{j\in\{2,4\}}
       \left(1-\frac{4F_LE_L}{j(j+2)}\right)
       \prod_{j\in\{2,4,\ldots,24\}}
       \left(1-\frac{4F_RE_R}{j(j+2)}\right).
\end{equation}
If \(\Pi t\) is a linear combination of \(\Pi n_i\), with
\(n_i\in N_4\), then \(t\in W_4\).
\end{lemma}
\begin{proof}
Both \(W_4\) and \(N_4\) are stable under current zero modes.
Every \(FE\) factor kills \([t]\in B_4/W_4\), so
\([\Pi t]=[t]\). Each \(\Pi n_i\) remains in \(N_4\subset W_4\).
Thus \([t]=0\). Equivalently, the weight-zero highest class is a singlet
in the finite semisimple current module. No dimension assertion about
the selected normal columns is needed.
\end{proof}

The prime search selects 21 independent projected normal columns.
All twelve projected identities have complete residual zero over
\(\QQ(\sqrt{13},\sqrt{-3})\). The final verifier repeats those
residuals by independent integer denominator clearing, authenticates
the raising premise, and repeats all nine saturation ranks. By
Lemmas~\ref{lem:raising} and~\ref{lem:singlet}, all \(9\cdot322=2898\)
required highest-left identities hold. Current covariance gives the
charged descendants. The second neutral field follows from
\(\omega_B-\omega_{K_L}-\omega_{K_R}\) minus the retained neutral field.
This proves \eqref{eq:gate4} with the certified lower controls.

\section{Gate 5 from Jacobi identities}\label{sec:gate5}
The proof uses the full normal spaces in
\eqref{eq:normal4}--\eqref{eq:normal5}, not only the selected solver
columns. Before using Gate 5, the lower controls and Gate 4 give
\begin{equation}\label{eq:normal-lemmas}
 (S_2)_0N_4\subset N_5,\qquad (N_3)_0S_3\subset N_5.
\end{equation}
The first follows by zero-mode derivation and skew reordering. For the
second, \((Ta)_0=0\) and, for a current \(x\),
\begin{equation}\label{eq:iterate}
 (x_{-1}a)_0p
 =\sum_{i\geq0}x_{-1-i}(a_ip)
  +\sum_{i\geq0}a_{-1-i}(x_ip).
\end{equation}
All inner products belong to the required lower spaces. To make the
translation reduction explicit, repeatedly use
\[
 x_{-r}v=\frac{T(x_{-(r-1)}v)-x_{-(r-1)}Tv}{r-1},\qquad r\geq2.
\]
The base terms are precisely those in \eqref{eq:normal5}; the same
identity applies to \(a_{-r}v\) with \(a\in S_2\). Thus no
unproved \(P_3\)-by-\(P_3\) closure enters \eqref{eq:normal-lemmas}.

If \(P_4\subset N_4\), Gate 5 already follows. Otherwise
\(Q=W_4/N_4\) is the irreducible current triplet of highest weight
\((0,2)\), since the actual \(P_4\) is irreducible. Let
\(q_-,q_0,q_+\) be its weight basis, and set
\[
 \beta(a,p)=[a_0p]\in Q,\qquad
 F(a,q)=[a_0\widetilde q]\in B_5/N_5.
\]
The second map does not depend on the representative
\(\widetilde q\in W_4\), by \eqref{eq:normal-lemmas}.

\begin{lemma}[Jacobi reduction]\label{lem:jacobi}
If \(a,b\in S_2\) satisfy \(a_0b\in N_3\), then for \(p\in S_3\),
\begin{equation}\label{eq:jacobi-reduction}
 F(a,\beta(b,p))=F(b,\beta(a,p)).
\end{equation}
\end{lemma}
\begin{proof}
The difference is the class of \((a_0b)_0p\), which vanishes by
\eqref{eq:normal-lemmas} and \([a_0,b_0]=(a_0b)_0\).
\end{proof}

\subsection{The charged families}
Let \(a_0,\ldots,a_7\) denote the eight charged highest fields, with
current weights
\[
 (0,4),(0,8),(0,8),(0,12),(1,1),(1,5),(1,7),(1,13).
\]
In a product such as \((a_i)_0p\), the parenthesized subscript denotes
the field label and the next subscript denotes its mode.
The saved exact relations give
\[
 \beta(a_3,p_3)=c_3q_-,\qquad \beta(a_6,p_6)=c_6q_-,
 \qquad c_3c_6\neq0,
\]
where \(p_3=\texttt{r3\_000910}\) has weight \((0,-14)\), and
\(p_6=\texttt{r3\_000899}\) has weight \((-1,-9)\).
The coefficients are authenticated in the final Jacobi certificate.
Exact combinations of the lower relations establish \((a_i)_0a_j\in N_3\)
for the following pairs:
\begin{center}
\begin{tabular}{@{}cccl@{}}
\toprule
\(i\)&\(j\)&Input&Weight of \((a_i)_0p_j\)\\
\midrule
0&3&\(p_3\)&\((0,-10)\)\\
1,2&3&\(p_3\)&\((0,-6)\)\\
3&6&\(p_6\)&\((-1,3)\)\\
4&6&\(p_6\)&\((0,-8)\)\\
5&3&\(p_3\)&\((1,-9)\)\\
6&3&\(p_3\)&\((1,-7)\)\\
7&3&\(p_3\)&\((1,-1)\)\\
\bottomrule
\end{tabular}
\end{center}
Every listed output weight lies outside \(Q\), so
\(\beta(a_i,p_j)=0\). Equation~\eqref{eq:jacobi-reduction} and
\(c_j\neq0\) give \(F(a_i,q_-)=0\). Raise this resulting zero
with \(a_i\) held highest to obtain the other two right inputs.
Current covariance gives all left descendants.

\begin{remark}
The weight exclusion is used only at the displayed bottom inputs.
It need not hold after raising \(p_j\); using it there would be an
incorrect stronger assertion. We raise the proved zero identity instead.
\end{remark}

\subsection{The neutral family}
Let \(C\) be the full current span, inside \(P=S_3/N_3\), of
brackets of charged weight-two fields. The charged-span certificate
reconstructs 28 actual charged seeds from the exact lower relations.
Together with the 644 lowering images, of rank 294, these have rank 322.
The independent verifier also checks both \(\mathfrak{sl}_2\) commutators
and all cross-axis commutators against the saved exact raising matrices.
Since \(C\) is a current submodule, the nilpotent-ideal argument gives
\[
 P=F_LP+F_RP+C\quad\Longrightarrow\quad P/C=0.
\]
Once again the finite-prime full rank is a positive spanning witness
inside an independently exact 322-dimensional space.

For the retained neutral singlet \(e\), write
\[
 e_0a_3=\sum_i t_i(x_i)_0y_i+n,
 \qquad n\in N_3,
\]
with charged \(x_i,y_i\). Apply zero modes to \(p_3\). By Jacobi,
the charged conclusion and \eqref{eq:normal-lemmas},
\[
 (e_0a_3)_0p_3
 =\sum_i t_i[(x_i)_0,(y_i)_0]p_3+n_0p_3\in N_5.
\]
Indeed Gate 4 puts each inner product in \(W_4\), and the next charged
product is controlled on both \(N_4\) and \(P_4\). A second Jacobi
identity gives
\[
 e_0((a_3)_0p_3)=(a_3)_0(e_0p_3)+(e_0a_3)_0p_3\in N_5.
\]
The nonzero coefficient of \(q_-\) proves the neutral bottom identity;
covariance gives the triplet. The other neutral field is obtained from
the total stress minus the two affine Sugawara stresses and \(e\).
This proves Gate 5. Its 24 independent highest channels comprise
23 charged channels and one neutral channel; they are consequences
of the displayed identities, not 24 separately recomputed grade-five
lattice products.

\section{From the finite gates to properness}\label{sec:properness}
The order of deductions prevents a circular use of Virasoro closure.
The direct controls in Certificate~\ref{cert:lower} first give
\(L(1)W_4\subset S_3\) and \(L(2)W_4\subset S_2\).
For quasiprimary \(a\in S_2\), Gate 4 then gives, for \(p\in P_3\),
\begin{align*}
 2a_1p&=L(1)(a_0p)-a_0L(1)p\in S_3,\\
 a_2p&=L(1)(a_1p)-a_1L(1)p\in S_2,\\
 2a_3p&=L(2)(a_1p)-a_1L(2)p\in K.
\end{align*}
The next mode is scalar. Affine normal left fields follow by iterate
identities. The actual provenance of \(P_4\) and
\[
 k_1(a_0p)=(k_0a)_1p+(k_1a)_0p+a_0(k_1p)
\]
give \(K_1P_4\subset S_3\). Current zero-mode stability is already
known; higher modes follow from the Virasoro commutators.
Consequently \(L(1)N_5\subset W_4\) and \(L(2)N_5\subset S_3\).
Gate 5 now yields, for \(p\in P_4\),
\begin{align*}
 2a_1p&=L(1)(a_0p)-a_0L(1)p\in W_4,\\
 a_2p&=L(1)(a_1p)-a_1L(1)p\in S_3,\\
 2a_3p&=L(2)(a_1p)-a_1L(2)p\in S_2,\\
 a_4p&=L(2)(a_2p)-a_2L(2)p\in K.
\end{align*}
The next mode is scalar. The 30 earlier prime-only \(a_1P_4\)
tests are therefore derived exactly here; they are not an omitted gate.

All hypotheses of Theorem~\ref{thm:recursive} hold for
\(K,S_2,S_3,P_4\). Their normal span is a vertex subalgebra of \(U\)
containing \(K,u,\omega_B\), hence equals \(U\). A normal word in
positive-weight generators cannot acquire weight one unless it is a
weight-one generator. Thus \(U_1=K\), proving
Proposition~\ref{prop:mainproper}.

\begin{reviewnote}{F1: what the finite certificate establishes}
This proof of properness has the strongest computational support in the
manuscript. It still relies on correctness of the lattice mode engine,
the saved actual-vector recipes, and the stated algebraic lemmas.
Hashes authenticate the records; hashes alone do not prove those lemmas
or independently reconstruct every earlier mode calculation. No claim
of a proof-assistant check or a second implementation of the entire
construction is intended. In particular, \(\dim U_4=6010\) is not used
as an input to this proof.
\end{reviewnote}

% Fragment: regularity, branching, fusion-ring reduction, and characters.
% The main document supplies theorem environments, reviewnote, and artifact.

\section{Regularity from a positive conformal-vector identity}
\label{sec:regularity}

This section concerns the actual PCT-stable subalgebra $U\subset B$ defined
in the construction. Its conclusions do not require properness. We use
the convention $Y(a,z)=\sum_{n\in\mathbb Z}a_nz^{-n-1}$ and write
$L^e(n)=e_{n+1}$ for the Virasoro operators of a conformal vector $e$.
The input from the local construction is the identity
\begin{equation}
 \omega_U=\omega_{A_{1,1}}
   +\frac{3403}{4563}\omega_{A_{1,39}}
   +\frac{25}{6}\overline p+\frac{25}{26}\overline z.
 \label{eq:positive-stress}
\end{equation}
Here $p$ and $z$ are the conformal vectors of actual embedded copies of
the three-state Potts algebra $K(\mathfrak{sl}_2,3)$ and the parafermion
algebra $K(\mathfrak{sl}_2,13)$; a bar denotes normalized averaging over
the compact current group. Each subalgebra is regular and unitary, and
its intrinsic grading agrees with the ambient grading on its states.
The averages take place in the finite-dimensional space $U_2$.
Finite-dimensional convexity replaces each average by a finite convex
combination of conjugate conformal vectors, so that
\begin{equation}
 \omega_U=\sum_{i=1}^N c_i e_i,\qquad c_i>0.
 \label{eq:finite-positive-stress}
\end{equation}
Zero coefficients, if any, are omitted. All simple modules of each
regular subalgebra $A_i$ have nonnegative lowest weights, with zero
occurring only for its vacuum module.

\begin{reviewnote}{R1: the embedded regular factors are local inputs}
The existence and normalization of the Potts and parafermion fields,
their inclusion in the generated $U$, and
\eqref{eq:positive-stress} are imported from the earlier construction
and the proof records
\artifact{hi_cft_20260923_next_structure_c2_energy_bound.md} and
\artifact{hi_cft_20260923_next_structure_top_bound.md}.
They are not consequences of the finite closure theorem alone.
This manuscript gives the deductions from those inputs, but does not
reproduce their full field-identification calculation. A review of
regularity must verify these embeddings, not only the displayed
positive coefficients.
\end{reviewnote}

\begin{lemma}[Orthogonal energy bound]
Let $V$ be a unitary VOA of CFT type admitting an identity of the form
\eqref{eq:finite-positive-stress}, with the properties just stated.
Then $V$ is $C_2$-cofinite.
\end{lemma}
\begin{proof}
For an ordinary simple $A_i$-module $M$, put
\[
 C_{2,A_i}(M)=\operatorname{span}\{a_{-2}m:a\in A_i,\ m\in M\}.
\]
These quotients are finite-dimensional by
\href{https://arxiv.org/html/math/0204021#S5}{Abe--Buhl--Dong,
Proposition~5.2}. There are finitely many simple $A_i$-modules. Hence
there is a finite number $D_i$ bounding the intrinsic energy in all
their $C_{2,A_i}$ quotients; this includes their lowest weights.
Regularity decomposes the weak $A_i$-module $V$ into ordinary simples,
and the unitary form permits an orthogonal decomposition. Since
\[
 C_{2,A_i}(V)\subset C_2(V),
\]
the self-adjoint intrinsic energy satisfies
\[
 \langle L^{e_i}(0)v,v\rangle\le D_i\|v\|^2
 \qquad (v\in C_2(V)^\perp).
\]
For $v\in V_n\cap C_2(V)^\perp$, the conformal-vector identity gives
\[
 n\|v\|^2
 =\sum_i c_i\langle L^{e_i}(0)v,v\rangle
 \le\Bigl(\sum_i c_iD_i\Bigr)\|v\|^2.
\]
Thus the graded quotient $V/C_2(V)$ vanishes above a finite weight.
No simultaneous diagonalization of the different energy operators is
required. The same proof applies directly to averages, since every
conjugate factor has the same bound $D_i$.
\end{proof}

\begin{lemma}[Lowest weights and projectivity of the vacuum]
Every simple ordinary $U$-module has nonnegative lowest conformal weight,
and the only simple module of lowest weight zero is $U$. Moreover $U$
is projective in the category of grading-restricted generalized modules.
\end{lemma}
\begin{proof}
Let $E$ be the lowest space of a simple ordinary module of weight $h$.
Every mode of an $A_i$ state that lowers intrinsic weight also lowers
ambient weight, and therefore annihilates $E$. Thus
$E\subset\Omega_{A_i}(M)$, the Zhu highest space. Restriction initially
gives a weak $A_i$-module; regularity decomposes it into ordinary simple
modules. The highest space of an ordinary simple is its top, so each
$L^{e_i}(0)|_E$ is diagonalizable with nonnegative real spectrum.
Taking traces, rather than assuming the operators commute, gives
\begin{equation}
 h\dim E=\sum_i c_i\operatorname{Tr}_E L^{e_i}(0)\ge0.
 \label{eq:top-trace}
\end{equation}
If $h=0$, every summand vanishes. Hence every partial energy is zero
on $E$, which lies in sums of vacuum tops for every $A_i$. Consequently
$L^{e_i}(-1)E=0$, and \eqref{eq:finite-positive-stress} gives
$L_U(-1)E=0$. Any nonzero $v\in E$ is vacuum-like: translation covariance
and lower truncation show that $Y_M(a,z)v$ has no negative powers of $z$.
The map $a\mapsto a_{-1}v$ is a nonzero module map $U\to M$, hence an
isomorphism.

Now let $\pi:E\twoheadrightarrow U$ be an epimorphism in the generalized
module category. $C_2$-cofiniteness gives finite length, and the preceding
argument excludes negative weights in all its composition factors.
The finite generalized weight-zero space $E_{[0]}$ therefore has no
lower-weight target for a lowering mode. It is again in every
$\Omega_{A_i}(E)$. On this space $L_U(0)$ is nilpotent, so
\[
 0=\operatorname{Tr}_{E_{[0]}}L_U(0)
  =\sum_i c_i\operatorname{Tr}_{E_{[0]}}L^{e_i}(0).
\]
Diagonalizability and nonnegativity of each partial energy imply it is
zero. The vacuum-top argument gives $L_U(-1)E_{[0]}=0$. Choose
$v\in E_{[0]}$ with $\pi(v)=\mathbf1$. Then
\[
 s:U\longrightarrow E,\qquad s(a)=a_{-1}v,
 \qquad \pi s=\mathrm{id}_U,
\]
is a section. Every epimorphism onto the vacuum splits.
\end{proof}

\begin{proposition}[Regularity of the actual generated algebra]
The VOA $U$ is strongly rational.
\end{proposition}
\begin{proof}
Simplicity, CFT type, and self-contragredience follow from the unitary
construction. The preceding lemmas supply $C_2$-cofiniteness and a
projective vacuum. The category $\mathcal C$ of grading-restricted
generalized modules is finite abelian and braided, with right-exact
tensor product and Grothendieck--Verdier duality whose dualizing object
is $U'=U$; the relevant VOA tensor-category hypotheses are reviewed in
\href{https://arxiv.org/html/2108.01898}{McRae, Sections~2.3 and~3}.

Here is the remaining categorical argument. Realize $\mathcal C$ as
finite-dimensional modules for a finite-dimensional algebra, denote
underlying vector-space dimension by $d_0$, and let $P$ correspond to
its regular module. Every $Z$ is a quotient of $P^{\oplus d_0(Z)}$.
For fixed $X$, right exactness gives
\[
 d_0(Z\otimes X)\le C_Xd_0(Z),\qquad
 \dim\operatorname{End}(X^{\otimes n})
 \le d_0(\mathbf1)^2C_X^{2n}<n!
\]
for large $n$, where $C_X=\max\{1,d_0(P\otimes X)\}$.
For nonzero indecomposable $X$, Grothendieck--Verdier evaluation
$DX\otimes X\to\mathbf1$ is nonzero, hence surjective onto the simple
unit, and splits by projectivity. Thus $X$ is non-negligible and has
moderate growth in the precise sense of
\href{https://arxiv.org/html/2412.17681v3}{Etingof--Penneys,
Theorem~1.1}; that theorem proves $X$ rigid. Finite direct sums give
rigidity of all objects. Tensoring is now exact, and
\[
 \operatorname{Hom}(X,-)
 \cong\operatorname{Hom}(\mathbf1,-\otimes X^*)
\]
is exact. Every object is projective, so $\mathcal C$ is semisimple and
$U$ is rational. This does not assume complete unitarity of its modules.
The local argument is recorded in
\artifact{hi_cft_20260923_next_structure_rationality.md}.
\end{proof}

\section{Branching and the boundary fusion ring}
\label{sec:branching}

From this point we use the properness conclusion and $\dim U_1=6$ of
the finite closure argument. Before that conclusion is supplied, every
statement in this section is conditional on the proper branch.

\subsection{A universal top bound and actual lattice branching}

The sharper bound uses an actual regular subalgebra $A_L\subset U$.
In the notation of the local construction, a rational lattice--Potts
core $\mathcal T$ and $K(\mathfrak{sl}_2,13)$ have pointed charges
$(j,k)\in(\mathbb Z/13)^2$ with quadratic form
$-(j^2+k^2)/13$. Exact weight-two vectors witness the simple-current
charge $(5,1)$ (up to reversing the second charge). Locality and
multiplicity one in a connected extension force its entire isotropic
line, yielding
\[
 A_L=\bigoplus_{k=-5j}\mathcal T_{1,j}\otimes\mathrm{PF}_{0,k}
 \subset U,\qquad
 \operatorname{Rep}(A_L)\simeq
 \mathrm{Fib}\boxtimes\mathrm{SO}(3)_{13}.
\]
The field witness and extension argument are in
\artifact{hi_cft_20260923_next_structure_order13_core.md} and
\artifact{hi_cft_20260923_next_structure_core_extension.json}.
There are fourteen simple tops, labelled $x_a=(1,a)$ and
$y_a=(\tau,a)$ for $0\le a\le6$.

Let $E$ be any simple $U$-module top and write $x_a,y_a$ also for the
nonnegative integer multiplicities of these $A_L$ tops in $E$.
Set $D=\dim E$. For each top, let $L$ be the trace of the left affine
energy, $Q=\operatorname{Tr}H_R(0)^2$, and $P,Z$ the Potts and
parafermion energy traces. Its contribution to \eqref{eq:top-trace} is
\begin{equation}
 W=L+\frac{83}{6084}Q+\frac{25}{6}P+\frac{25}{26}Z.
 \label{eq:top-W}
\end{equation}
This is a trace formula on a complete right $\mathfrak{sl}_2$
representation, not an assertion about a single partial-energy
eigenvector. The exact dimensions and trace contributions are
\[
\begin{array}{c|r|r@{\qquad}c|r|r}
 \text{top}&d&4680W&\text{top}&d&4680W\\\hline
 x_0&1&0&y_0&1&7800\\
 x_1&1&600&y_1&7&36960\\
 x_2&5&10440&y_2&15&110880\\
 x_3&9&56640&y_3&2&4080\\
 x_4&2&5520&y_4&14&81480\\
 x_5&14&148560&y_5&6&18240\\
 x_6&6&54240&y_6&2&4920
\end{array}
\]
Complete right $\mathfrak{sl}_2$ characters imply
$y_2\ge x_5$ and $y_4\ge x_5+x_6$: these compare charges $10,12$
and charges $7,9$, respectively, with the common left doublet removed
in the second comparison. Direct comparison of the displayed table gives
\begin{align}
 4680\bigl(\tfrac{17}{10}D-W_{\rm total}\bigr)
 ={}&8460(y_2-x_5)+28716(y_4-x_5-x_6)\notag\\
 &+7956x_0+7356x_1+29340x_2+14964x_3+10392x_4\notag\\
 &+22212x_6+156y_0+18732y_1+11832y_3\notag\\
 &+1188y_4+29496y_5+10992y_6\ge0.
 \label{eq:sharp-top-bound}
\end{align}
Since $W_{\rm total}=hD$, this proves $0\le h\le17/10<2$ for
\emph{every} simple $U$-module, without assuming it is unitarizable.

\begin{reviewnote}{R2: the core-top table and finite-top actions}
The inequality \eqref{eq:sharp-top-bound} can be checked directly from
the displayed table. Deriving that table, including all minimizing
charge representatives, uses
\artifact{hi_cft_20260923_next_structure_order13_top_bound.md}.
The construction of $\mathcal T$, the thirteen-charge extension, and
completeness of the resulting fourteen-module list are substantial
local inputs whose full derivations are not repeated here. Likewise,
the finite-top irreducibility used below is imported from exact field
actions, particularly
\artifact{hi_cft_20260923_next_81top_all.json}; it should not be replaced
by numerical cyclicity alone.
\end{reviewnote}

The strongly rational restriction $B\downarrow_U$ is semisimple. Its
simple summands have integer lowest weights, hence only weights $0,1$
are possible. The vacuum occurs once, while the complementary current
space has dimension $86-6=80$. Its certified irreducible zero-mode
action therefore implies
\begin{equation}
 B\downarrow_U=U\oplus X,\qquad h_X=1,\qquad\dim X_{\rm top}=80.
 \label{eq:B-branching}
\end{equation}
The four mixed nontrivial $E_6\oplus A_2$ lattice sectors start at
weight $1$, with $81$-dimensional tops. Their next possible top weight
would be $2$, excluded by \eqref{eq:sharp-top-bound}. Their exact
top-action certificates use a one-dimensional Cartan eigenline,
positive-metric adjoint identities, and a cyclic span of dimension
$81$. The adjoint identities make the generated operator algebra a
star algebra; orthogonal complements are therefore invariant, and
the unique eigenline plus cyclicity proves irreducibility.

Write the actual lattice extension as $H=B\oplus J\oplus J^2$. Its
restriction is consequently
\begin{equation}
 H\downarrow_U=\mathbf1\oplus X\oplus Y\oplus Y^*,\qquad
 h_Y=1,\quad\dim Y_{\rm top}=81.
 \label{eq:H-branching-fork}
\end{equation}
At this stage $Y\cong Y^*$ has not been assumed. The three types
$\mathbf1,X,Y$ are distinguished by their lowest spaces.

\subsection{The rank-four or rank-six alternative}

Put $\mathcal C=\operatorname{Rep}(U)$ and
$\mathcal F=\mathcal C_H$, the category of \emph{all} right modules
over the extension algebra $H$. The strict positivity of nonvacuum
lowest weights gives the positive spherical structure through character
asymptotics. Since $H$ is holomorphic, it is Lagrangian and the
\href{https://arxiv.org/abs/1009.2117}{DMNO center correspondence}
gives $\mathcal C\simeq Z(\mathcal F)$.
Write $s=\operatorname{FPdim}_{\mathcal C}B$. Extension dimension
identities give
\[
 d_X=s-1,\qquad d_Y=d_{Y^*}=s,\qquad
 \operatorname{Dim}\mathcal C=9s^2,\quad
 \operatorname{Dim}\mathcal F=3s.
\]
Pseudo-unitarity of $\mathcal F$ follows from that of its center and
the squared global-dimension and Frobenius--Perron-dimension identities.
No $C^*$ structure is needed in this section.

For center induction $I$ and forgetful functor $F$, one has $H=I(\mathbf1)$
and $FI(\mathbf1)=\bigoplus_{a\in\operatorname{Irr}\mathcal F}a\otimes a^*$.
Adjunction therefore gives
\begin{equation}
 \operatorname{rank}\mathcal F
 =\dim\operatorname{End}_{\mathcal C}(H)
 =\begin{cases}4,&Y\not\cong Y^*,\\6,&Y\cong Y^*.
 \end{cases}
 \label{eq:rank-fork}
\end{equation}
The actual order-three automorphisms of $H/B$ inject into
$\operatorname{Inv}(\mathcal F)$. This injection is already spherical:
under the composition/convolution description of
$\operatorname{End}(I(\mathbf1))$ in
\href{https://arxiv.org/html/1909.02816}{Bischoff--Jones,
Lemma~3.1 and Corollary~3.4}, an algebra automorphism and its inverse
are convolution idempotents. They are positive sums of the primitive
idempotents $d_a[a]$; their fusion product is $[\mathbf1]$. Positivity
forces singleton supports consisting of inverse invertibles.

The earlier pointwise automorphism calculation gives
$\operatorname{Aut}(B/U)=1$. If rank six had six invertibles, then
$3s=6$, so $X$ would be invertible and $B=\mathbf1\oplus X$ would
have a nontrivial sign automorphism, a contradiction. Thus there are
exactly three invertibles in either rank. Their associator is also
determined: the local $B$-module category has metric group
\[
 (\mathbb Z/3)^2,\qquad q(a,b)=\frac{2a^2+b^2}{3}\pmod1.
\]
The actual extension $H/B$ selects an isotropic line. Forming its
modules gives an embedded untwisted $\operatorname{Vec}_{\mathbb Z/3}$
subcategory of $\mathcal F$. Merely knowing the group of invertibles
would not determine this associator.

\subsection{Excluding rank four: scope of the longer local argument}

At rank four the fusion ring is near-group:
\[
 a^3=1,\qquad a\rho=\rho a=\rho,\qquad
 \rho^2=1+a+a^2+m\rho.
\]
The published \href{https://doi.org/10.1016/j.jalgebra.2022.02.020}
{Dong--Chen--Wang, Theorem~6.2}, restricts $m$ to $0,2,3,6$.
For $m=0$, $s=2$ and the sign-automorphism contradiction just given
applies. For $m=2$, the categorical reduction of
\href{https://arxiv.org/html/math/0304194#S7}{Etingof--Gelaki--Ostrik,
Theorem~7.1 and Corollary~7.4} supplies the twisted $A_4$ double.
Its electric regular algebra gives a holomorphic $c=8$ extension
$H'$ with $U=(H')^{A_4}$, and
\href{https://msp.org/pjm/2004/213-2/pjm-v213-n2-p05-p.pdf}
{Dong--Mason, Theorem~1}, identifies $H'$ with $V_{E_8}$.
The order-two and order-three adjoint traces in $E_8$ are bounded
below by $-8$ and $-4$, respectively, as follows from the affine
Dynkin-diagram classification. Therefore
\[
 \dim(E_8^{A_4})\ge\frac{248+3(-8)+8(-4)}{12}=16,
\]
contradicting $\dim U_1=6$.

For clarity, the remaining exclusions use universal spherical
induction arguments, not a classification limited to unitary examples.
For $m=3k$, induction yields integers $p,q\ge0$ with $p+q=k$.
At $k=1$ the pointed factor and two indicator moments force the modular
representation of $G_{2,3}\otimes A_{2,1}$ or its reversal. Harmonic
theta-series pole cancellation in these two orientations gives
\[
 \dim U_1=15+A+10B\quad\hbox{or}\quad14+42A,
 \qquad A,B\in\mathbb Z_{\ge0}.
\]
At $k=2$, the balanced pattern $p=q=1$ forces, by signed Galois symmetry,
a row of the center's $S$ matrix whose Verlinde coefficient is
\[
 N_{XX}^{X}=\frac{97}{72}+\frac{32}{72}-\frac{33}{72}=\frac43.
\]
The unbalanced pattern $\{p,q\}=\{0,2\}$ admits a weighted character
projection. Its three-dimensional block is
\[
 S_v=-\frac12\begin{pmatrix}1&2&1\\1&0&-1\\1&-2&1\end{pmatrix},
 \qquad T_v=\operatorname{diag}(1,-i,-1)
\]
before the common central-charge factor; reversal conjugates $T_v$.
Harmonic theta-series subtraction gives the two alternatives
\[
 \dim U_1=8+A\quad\hbox{or}\quad5+12P,
 \qquad A,P\in\mathbb Z_{\ge0}.
\]
None of these formulas allows six currents. The parameters are actual
weighted polar top dimensions, so their integrality and positivity
do not require linear independence of all simple characters.

\begin{reviewnote}{R3: compressed rank-four exclusions}
The preceding paragraph is a summary, not a full proof of the universal
$m=3,6$ exclusions. The derivation of the complete modular representation
for $m=3$ is in
\artifact{hi_cft_20260923_next_m3_universal_data.md}, with an independent
audit in
\artifact{hi_cft_20260923_next_m3_universal_independent_audit.md}.
The $m=6$ arguments are in
\artifact{hi_cft_20260923_next_structure_m6_balanced_exclusion.md} and
\artifact{hi_cft_20260923_next_structure_m6_universal_projection.md}.
The coefficient calculations and pole-cancellation arguments are in
\artifact{hi_cft_20260923_next_structure_m3_character_obstruction.md}
and
\artifact{hi_cft_20260923_next_structure_m6_character_obstruction.md}.
These are new local deductions, not statements of the cited
rank-four classification theorem. Their universal induction reductions
and theta normalizations remain an explicit review dependency of this
condensed manuscript. In particular, the older character notes alone
only exclude specified model representations; the later universal
reductions are essential. The pointwise claim
$\operatorname{Aut}(B/U)=1$ and the $m=2$ electric-extension argument
also rely on the longer record
\artifact{hi_cft_20260923_next_structure_rank46.md} and its replacement
spherical audit
\artifact{hi_cft_20260923_next_structure_final_dependency_audit.md}.
\end{reviewnote}

\subsection{The rank-six ring and exclusion of higher multiplicity}

Rank four being excluded, \eqref{eq:rank-fork} gives
$H=\mathbf1\oplus X\oplus2Y$. The fusion algebra is
\[
 K_0(\mathcal F)\otimes\mathbb C
 \cong\mathbb C\oplus\mathbb C\oplus M_2(\mathbb C),
\]
by the composition algebra of $\operatorname{End}_{\mathcal C}(H)$.
Hence it is noncommutative. The three invertibles act transitively on
the three noninvertibles: otherwise both nontrivial idempotents of
$\mathbb C[\mathbb Z/3]$ split off, leaving a four-dimensional ideal
with a dimension character, incompatible with the $M_2$ block.
The left and right cyclic actions must be inverse; equality would
make the fusion ring commutative. Duality then fixes every noninvertible,
and reciprocity and associativity give
\begin{equation}
 a^3=1,\qquad\rho a=a^2\rho,\qquad
 \rho^2=1+m(\rho+a\rho+a^2\rho),\qquad m\ge1.
 \label{eq:rank6-ring}
\end{equation}

Here are the essential numerical identities in the local multiplicity
exclusion. Put
\[
 d=\frac{3m+\sqrt{9m^2+4}}2,\quad s=1+d^2,\quad
 \Delta=9m^2+4,\quad R=\rho+a\rho+a^2\rho.
\]
Center induction, with the untwisted pointed subcategory, yields
\[
 I(1)=1+X+2Y,\quad d_X=d^2,\quad d_Y=s,
 \qquad I(a)=I(a^2)=C_0+C_1+C_2,
\]
where $\theta_X=\theta_Y=1$, $d_{C_j}=s$, and the $C_j$ twists are
$1,\zeta_3,\zeta_3^2$.
All remaining central simples $Z_j$ have $F(Z_j)=n_jR$,
$n_j\in\mathbb Z_{>0}$. Global dimension and the first two
Ng--Schauenburg indicators give
\begin{align}
 \sum_jn_j^2&=4m^2+2,\notag\\
 \sum_jn_j^2\theta_j&=-\frac{m^2+m\sqrt\Delta}{2},
 \label{eq:indicator-first}\\
 \sum_jn_j^2\theta_j^2
 &=\frac{-m^2+(2\varepsilon-m)\sqrt\Delta}{2},
 \qquad\varepsilon=\nu_2(\rho)\in\{1,-1\}.
 \label{eq:indicator-second}
\end{align}
These follow from the center-twist trace formula in
\href{https://arxiv.org/pdf/math/0601012}{Ng--Schauenburg,
Theorem~4.1}. They concern a multiset of exactly $4m^2+2$ roots,
repeating $\theta_j$ precisely $n_j^2$ times.

Write $\Delta=t^2D$ with $D$ squarefree. One has $D\equiv1\pmod3$
and $7\nmid D$. Except when $D=10$, the elementary prime-factor estimate
$\varphi(2D)/\sqrt D\ge12/\sqrt{13}$ and
\href{https://arxiv.org/pdf/1401.1879}{Larson's root-count theorem}
applied to \eqref{eq:indicator-first} give
\[
 4m^2+2\ge\frac{mt}{2}\varphi(2D)
 \ge\frac{6m\sqrt{9m^2+4}}{\sqrt{13}}.
\]
For $m>1$ this is impossible, since
\[
 36m^2(9m^2+4)-13(4m^2+2)^2
 =4(m^2-1)(29m^2+13)>0.
\]
For $D=10$, a root multiset with sum $-a-b\sqrt{10}$, with
$a,b\in\mathbb Z_{\ge0}$, has at least $a+8b$ members. In the equality
case its squared-root sum is $a$, provided that sum lies in
$\mathbb Q(\sqrt{10})$. To see the equality condition, average over
the Galois group fixing $\sqrt{10}$. Rational contributions cost at
least one root per negative unit, with equality only for $-1$;
quadratic contributions cost at least eight roots per unit, with
equality only for a primitive-40 orbit. That orbit has squared sum
zero. The hypothesis on the second sum permits undoing the average.

In our application $m,t$ are even and $a=m^2/2$, $b=mt/2$.
The bound
\[
 4m^2+2\ge\frac{m^2}{2}+4mt
\]
excludes even $m\ge4$. The remaining $m=t=2$ is the equality case,
so the squared sum must be $2$, whereas \eqref{eq:indicator-second}
gives $-2$ or $-2-4\sqrt{10}$. Thus $m=1$.
The full conductor-orbit proof, including imprimitive orbits and the
second-moment hypothesis, is retained in
\artifact{hi_cft_20260923_next_structure_rank6_indicators.md} and
\artifact{hi_cft_20260923_next_rank6_root_audit.md}.

Consequently
\begin{equation}
 d=\frac{3+\sqrt{13}}2,\qquad
 s=\frac{13+3\sqrt{13}}2,\qquad
 d_X=\frac{11+3\sqrt{13}}2.
 \label{eq:haagerup-dimensions}
\end{equation}
For $m=1$, \eqref{eq:indicator-first} has six roots with sum
$(-1-\sqrt{13})/2$. The quadratic-period equality case forces the
six distinct nonresidue thirteenth roots, so $n_j=1$ and the center
has twelve simples. Its dimensions and twist exponents are
\[
 \begin{array}{c|cccc}
 \text{number}&1&1&4&6\\\hline
 \text{dimension}&1&d^2&s&3d
 \end{array}
 \qquad
 0,0,0,0,\tfrac13,\tfrac23,
 \tfrac2{13},\tfrac5{13},\tfrac6{13},
 \tfrac7{13},\tfrac8{13},\tfrac{11}{13}.
\]
\begin{reviewnote}{R4: conductor-orbit bounds and equality cases}
The reduction to $m=1$ and the count of twelve central simples use more
than the displayed indicator identities. The prime-factor estimate,
the exceptional $\mathbb Q(\sqrt{10})$ equality case, and the
quadratic-period equality case at $m=1$ are compressed here. Their full
conductor-orbit arguments, including imprimitive roots and multiplicities,
are in
\artifact{hi_cft_20260923_next_structure_rank6_indicators.md},
\artifact{hi_cft_20260923_next_rank6_root_audit.md}, and
\artifact{hi_cft_20260923_next_structure_rank6_center.md}.
These local deductions need review independently of the cited indicator
formula. The later net-sector count inherits this dependency.
\end{reviewnote}
The complete $S,T$ data follow from
\href{https://arxiv.org/html/2609.25185v1}{Gannon--Schopieray--Yadav,
Theorem~A(2), equivalently Theorem~5.1}: its relevant statement assumes
a complex spherical untwisted HI category. The auxiliary metric group
has prime order $13$ and its negative quadratic Gauss sum fixes its
isometry class. This identifies the standard Haagerup modular data.
It does not by itself identify a braided category; that stronger
conclusion is treated separately after complete unitarity.

\section{The complete character vector}
\label{sec:characters}

Let $q=e^{2\pi i\tau}$. Character exponents are $h-c/24=h-1/3$.
Using the preceding twist list and strict positivity of nonvacuum
lowest weights, the only allowed nonvacuum pole is
$\gamma q^{-7/39}$ in the twist-$2/13$ coordinate. Here
$\gamma\in\mathbb Z_{\ge0}$ is its top dimension if that module starts
at $2/13$, and is zero if its lowest weight is shifted upward. The
vacuum pole is $q^{-1/3}$ with coefficient one.

\begin{lemma}[Uniqueness from the principal part]
Two character vectors for this modular representation with the same
$\gamma$ coincide in every weight.
\end{lemma}
\begin{proof}
Their difference is holomorphic at infinity. The modular representation
has finite image, and covariance transports this property to every cusp
of a finite-index kernel. Each component is therefore a holomorphic
function on a compact modular curve, hence constant. A constant vector
must be fixed by both $S$ and the \emph{character} $T$ matrix, which
includes $e^{-2\pi i/3}$. The only $T$-fixed coordinate is the object
of categorical twist $e^{2\pi i/3}$. Its $S$ image has a nonzero vacuum
component, so it cannot span an invariant line. The constant is zero.
\end{proof}

\href{https://arxiv.org/pdf/1006.1326}{Evans--Gannon, Sections~1.4
and~2.4, equation~(2.7)}, provide a fundamental matrix $\Xi$ whose
vacuum and $d_3$ columns have precisely the two principal parts above.
Thus, in their labeling,
\[
 \boldsymbol\chi=\Xi_{:,0}+\gamma\Xi_{:,d_3},
 \qquad\dim U_1=6+13\gamma.
\]
The actual six-current result forces $\gamma=0$. Their recursive
construction of $\Xi$ and the uniqueness lemma give the entire vector,
not just a finite coefficient match. In particular,
\begin{equation}
 \operatorname{ch}U
 =q^{-1/3}\bigl(1+6q+120q^2+956q^3+6010q^4+\cdots\bigr).
 \label{eq:vacuum-character}
\end{equation}
The lowest weights and top dimensions are
\[
\begin{array}{c|c|r@{\qquad}c|c|r}
 \text{label}&h&\dim M_{\rm top}&\text{label}&h&\dim M_{\rm top}\\\hline
 0&0&1&d_1&6/13&7\\
 b&1&80&d_2&11/13&42\\
 a&1&81&d_3&15/13&119\\
 c_0&1&81&d_4&5/13&5\\
 c_1&1/3&3&d_5&7/13&13\\
 c_2&2/3&27&d_6&8/13&14
\end{array}
\]
In particular the final maximum is $15/13$, improving the preliminary
$17/10$ bound. The proof record is
\artifact{hi_cft_20260923_next_structure_conditional_characters.md}.
Although this vector fixes every unflavoured homogeneous dimension,
it does not determine the separate branching functions for the
fourteen-module core $A_L$, and it was not used to prove properness.

% Fragment: analytic comparison, unitary identification, and full theories.
% No preamble.  The master document supplies theorem-like environments,
% criticalnote/reviewnote, \artifact, \Rep, \Com, \FPdim, \Vec, and \HH.

\section{From vacuum unitarity to complete unitarity}
\label{sec:analytic-unitarity}

The exact strong-generation certificates concern the actual subVOA
$U\subset B\subset H=V_{E_8}$.  In particular, they establish
$\dim U_1=6$ and $U\ne B$; they do not establish positivity on all
modules or on fusion products.  We now separate those questions.
\emph{Vacuum unitarity} is the positive invariant form inherited by the
PCT-stable subVOA $U$.  \emph{Strong unitarity} means that all simple
ordinary $U$-modules admit compatible positive invariant forms.
\emph{Complete unitarity} additionally includes positivity of the
canonical Hermitian fusion forms, giving the VOA representation category
its unitary modular tensor structure.  An abstract ribbon equivalence to
a unitary category does not, by itself, prove this last positivity statement.

We use the preceding algebraic analysis as input: regularity of $U$,
the actual restrictions
\begin{equation}
 B\mathord\downarrow_U=U\oplus X,
 \qquad
 H\mathord\downarrow_U=U\oplus X\oplus Y^{\oplus2},
 \qquad h_X=h_Y=1,
 \label{eq:analytic-branching}
\end{equation}
the twelve simple algebraic $U$-modules, and the
$\operatorname{HI}(\mathbb Z/3)$ fusion ring of the all-module category
for the original extension.  Every use of these inputs inherits the
qualifications attached to their preceding proofs.  Agreement of modular
data alone is never used as a classification of tensor categories.

\begin{criticalnote}{Analytic proof boundary}
The local applications labelled A1--A3 below are \textbf{claims pending
expert review}.  Their proposed arguments are included so that the
domains, geometric identifications, and counting steps can be checked.
The cited external theorems are used with their stated hypotheses;
the issue under review is whether our construction satisfies and matches
those hypotheses.  Thus complete unitarity, the unitary ribbon
identification, and the specifically $\HH_3$ extension remain conditional
in this manuscript.  The exact finite certificates are separate accepted
inputs and are not weakened by this distinction.
\end{criticalnote}

\subsection{The subnet and its sector count}
\label{subsec:analytic-net-count}

Write $\mathcal N_U\subset\mathcal N_H$ for the actual conformal-net
inclusion.  The lattice VOA supplies strong locality and bounded
localized vertex operators; the unitary conformal subalgebra inherits
the relevant local construction.  Once regularity is available,
\href{https://link.springer.com/article/10.1007/s00039-024-00685-8}
{Tener, Corollary~4.24}, applies to the conformal inclusion and gives
finite index.  The holomorphic lattice net is completely rational;
\href{https://arxiv.org/html/math/0102196}
{Longo, Theorem~24}, gives complete rationality of the subnet.
This finite-index argument does not require every $U$-module to be unitary.

\begin{criticalnote}{A2: actual net branching and twelve sectors}
\begin{claim}[Net-sector count; pending expert review]
\label{clm:analytic-net-count}
The restriction of the holomorphic $H$ net has precisely the three
inequivalent irreducible $\mathcal N_U$ summands corresponding to
\eqref{eq:analytic-branching}, with multiplicities $1,1,2$.
Its boundary fusion category $\mathcal F_{\mathrm{net}}$ has rank six,
untwisted invertible subcategory $\Vec_{\mathbb Z/3}$, and the
$\operatorname{HI}(\mathbb Z/3)$ fusion ring.  Consequently
\[
 \Rep(\mathcal N_U)\simeq Z(\mathcal F_{\mathrm{net}}),
 \qquad
 \#\operatorname{Irr}\Rep(\mathcal N_U)=12.
\]
\end{claim}
The passage from algebraic summands to irreducible net summands, and
the transfer of the actual automorphism and endomorphism information,
must be verified at the level of closed local operators.  A shared
branching matrix does not identify two tensor categories.
\end{criticalnote}

Here is the proposed reduction of this claim.  On finite-energy vectors
of $H$, reconstructed $U$-fields are its original restricted fields.
Orthogonal projections onto its actual simple summands should commute
with their closures.  The resulting Lagrangian algebra has endomorphism
algebra
\begin{equation}
 \operatorname{End}(A_{\mathrm{net}})
 \cong\mathbb C\oplus\mathbb C\oplus M_2(\mathbb C).
 \label{eq:analytic-end-algebra}
\end{equation}
The center-induction description then gives rank six and a
noncommutative complex fusion algebra for $\mathcal F_{\mathrm{net}}$.
The actual order-three automorphisms of $H/B$ supply three invertibles.
If all six simples were invertible, the intermediate index would be
$[\mathcal N_B:\mathcal N_U]=2$.  The connected intermediate algebra
$1+X$ would have a sign automorphism.  Reconstruction on finite energies
would turn it into an automorphism of $B$ fixing $U$, contrary to the
actual calculation $\operatorname{Aut}(B/U)=1$.  Thus the invertible
group is exactly $\mathbb Z/3$.  The pointed nine-sector $B$ net and the
actual $H/B$ extension supply its untwisted associator class.

With these positive spherical inputs, the ring and indicator arguments
apply independently on the net side.  In particular, for
$d=(3+\sqrt{13})/2$ and $s=1+d^2$, the center calculation gives dimensions
\begin{equation}
 1,\ d^2,\ \underbrace{s,s,s,s}_{4},\
 \underbrace{3d,3d,3d,3d,3d,3d}_{6}.
 \label{eq:analytic-center-dimensions}
\end{equation}
The last six twists have exponents
$\{2,5,6,7,8,11\}/13$.  This is the proposed source of the net count
twelve.  It is not obtained by assigning the VOA character transformation
matrix to the net category.  The underlying deductions are recorded in
\artifact{hi_cft_20260923_next_structure_rank46.md},
\artifact{hi_cft_20260923_next_structure_rank6_indicators.md}, and
\artifact{hi_cft_20260923_next_structure_rank6_center.md}.

\subsection{Reverse reconstruction and exhaustion}
\label{subsec:analytic-reconstruction}

\href{https://arxiv.org/html/2507.20735v3#S9}
{Henriques--Tener, Theorem~9.5 and Corollary~9.7}, reconstruct unitary
modules and identify their closed smeared fields.  Remark~9.8 extends
the construction beyond the finite-multiplicity hypothesis.  This
extension matters: finite-dimensional energy eigenspaces must not be
silently assumed for the net representation one is trying to reconstruct.

\begin{criticalnote}{A3: ordinaryness, irreducibility, and linear fullness}
\begin{claim}[Reverse module comparison; pending expert review]
\label{clm:analytic-reconstruction}
For this regular $U$, reconstruction sends each irreducible
$\mathcal N_U$ sector to an irreducible ordinary unitary $U$-module.
It preserves $L_0$ and induces a fully faithful linear $*$-functor on
finite direct sums,
\[
 \mathcal R:\Rep(\mathcal N_U)\longrightarrow\Rep^{u}(U).
\]
No tensorator or compatibility with braiding is asserted here.
\end{claim}
The projection argument below must establish common-core invariance and
commutation with closed local fields, including their adjoints.  These
are our analytic deductions from reconstruction, not a quoted
braided-equivalence theorem.
\end{criticalnote}

Let $\mathcal K$ be an irreducible net representation and let
\[
 M=\bigcup_{R>0}1_{[0,R]}(L_0)\mathcal K
\]
be the dense bounded-energy space appearing in the extended construction.
Regularity decomposes this admissible weak $U$-module algebraically into
ordinary simple modules.  Choose a nonzero summand $S$, inherit its
positive form, and let $P$ be the Hilbert projection onto
$\overline S\subset\mathcal K$.

The intended argument first uses the agreement of the two Virasoro
actions, as in Henriques--Tener Corollary~9.6.  Each energy space of $S$
is finite-dimensional, so its closure adds no vector at a fixed energy.
It follows that $P$ preserves the bounded-energy domain and maps it
into $S$.  The unitary adjoint identities, tested against $S$, then give
commutation with all modes on the common domain.  The required closure
argument is that $P$ and $1-P$ reduce the closed smeared fields and
therefore the local algebras they generate.  Irreducibility of the net
sector implies $P=1$, and the finite-dimensional energy spaces then
give $M=S$.  This proves ordinaryness and simplicity if the stated
domain assertions hold.

For fullness, an algebraic isomorphism between simple unitary modules is
unitary after multiplication by a scalar, by uniqueness of the positive
invariant form.  Its bounded extension should intertwine the closed
fields by the same argument.  Conversely, net intertwiners commute with
the reconstructed modes.  Semisimplicity handles finite direct sums.
The detailed local source is
\artifact{hi_cft_20260923_next_m3_unitarity_scope.md}.

\begin{proposition}[Conditional strong unitarity]
\label{prop:analytic-strong-unitarity}
Assume Claims~\ref{clm:analytic-net-count}
and~\ref{clm:analytic-reconstruction}, and the preceding algebraic
count of twelve simple $U$-modules.  Then every simple $U$-module is
unitarizable, and its character is the character of a net sector.
\end{proposition}
\begin{proof}
Full faithfulness makes the twelve net simples pairwise inequivalent
simple ordinary $U$-modules.  There are exactly twelve algebraic simple
modules, so this injection is a bijection.  Each image carries its
reconstructed positive form, and $L_0$ is preserved.  Finite direct sums
give the corresponding statement for ordinary modules.
\end{proof}

This counting step supplies essential surjectivity only as a linear
comparison.  Neither fusion compatibility, equality of categorical
$S$ matrices, nor positivity of VOA fusion forms follows from it.

\subsection{The annulus compatibility claim}
\label{subsec:analytic-annuli}

The next step concerns the specific geometric annulus system of
\href{https://arxiv.org/html/1810.08168#S3.SS4}
{Tener, Definition~3.18 and Corollaries~3.20--3.22}.
Let $(B,A)$ be a localized pair for an interval $I$, with incoming
factor $A$ and outgoing factor $B$.  Subscripts distinguish their
actions on the vacuum and on a reconstructed module.  For an insertion
point $z\in\operatorname{int}(B,A)$, put
\[
 X_v(B,z,A)=\overline{B_UY_U(v,z)A_U}\in\mathcal N_U(I).
\]
The bar denotes the bounded extension when it exists.  The compatibility
required by Tener's definition of a module net representation is the
specific operator identity in the following claim.

\begin{criticalnote}{A1: new annulus compatibility lemma}
\begin{claim}[Annulus compatibility; pending expert review]
\label{clm:analytic-annulus}
For the actual geometric system above, choose matching paths and
central lifts in the two annulus constructions.  If $M$ is an ordinary
unitary module reconstructed from a sector $\pi_M$, then
\begin{equation}
 \pi_M\bigl(X_v(B,z,A)\bigr)
   =\overline{B_MY_M(v,z)A_M},
 \qquad v\in U,\quad z\in\operatorname{int}(B,A).
 \label{eq:analytic-annulus-identity}
\end{equation}
The right side exists as a bounded extension.  The same matching in the
vacuum identifies the net's geometric insertion with $X_v(B,z,A)$.
\end{claim}
The points requiring expert verification are modulewise matching of
the annulus systems and central lifts, density and core properties of
the inverse-finite-energy domains, and holomorphic continuation on every
component across a possible circle seam.  Agreement of vacuum operators
alone would not prove this claim.
\end{criticalnote}

We give the proposed proof in three steps.  First, geometric localization
must precede operator comparison.  The thin-boundary parametrizations in
the selected system agree outside $I$, and its common localizer is
geometric.  Its generating operators are univalent-semigroup factors
$\exp(-tL(\rho))$ and positive-energy diffeomorphism or M\"obius
operators.  The time-ordered integration in
\href{https://arxiv.org/html/2506.08684#S5}
{Henriques--Tener, formula~(5.2)}, gives the proposed comparison with the
geometric annulus action.  Constant generators give the same exponential;
thin factors must give the same diffeomorphism action.  A common path
and central lift are essential.  Using the same localizing lift for both
members of $(B,A)$ is intended to cancel deck-transformation factors.
This matching is required on every positive-energy module, not merely
on the vacuum.

Second, let $M^{\mathrm{fin}}$ be the finite-energy vectors and define
\begin{equation}
 \mathcal D_A=\{\xi\in\mathcal K_M:A_M\xi\in M^{\mathrm{fin}}\},
 \qquad
 \mathcal D_{B^*}=\{\eta\in\mathcal K_M:B_M^*\eta\in M^{\mathrm{fin}}\}.
 \label{eq:analytic-annulus-domains}
\end{equation}
For an incoming annulus, choose a small round circle inside its smooth
inner Jordan domain, containing the origin, and fill the complementary
annulus by $C$.  The proposed factorization, with the matched central
scalar retained, is
\begin{equation}
 A_MC_M=\lambda r^{L_0},\qquad 0<r<1,\quad\lambda\ne0.
 \label{eq:analytic-round-factorization}
\end{equation}
It puts finite-energy vectors in $\operatorname{Ran}(A_M)$ and gives
$C_M M^{\mathrm{fin}}\subset\mathcal D_A$.  The dense-range property of
$C_M$ would make $\mathcal D_A$ dense; the adjoint construction treats
$\mathcal D_{B^*}$.  This is the domain issue addressed by the proposed
use of Henriques--Tener Lemma~7.9.

Third, Theorem~9.5(i), together with insertion gluing in
Proposition~6.19, supplies the incoming identity after the preceding
geometric matching.  Adjoints supply the outgoing identity.  For
$\xi\in\mathcal D_A$ and $\eta\in\mathcal D_{B^*}$, the algebraic pairing
\begin{equation}
 \left\langle Y_M(v,z)A_M\xi,B_M^*\eta\right\rangle
 \label{eq:analytic-laurent-pairing}
\end{equation}
is a finite Laurent polynomial: both vectors have finite energy support,
and the mode grading permits only finitely many powers.  The geometric
insertion pairing is holomorphic on $\operatorname{int}(B,A)$.
The proposed geometric argument shows that every connected component
meets an incoming or outgoing open part.  The identity theorem would
then extend equality over the component, including across its seam.
Equality on the two dense test domains identifies the bounded extension
with $\pi_M(X_v(B,z,A))$.  Applying the same reasoning to the vacuum
finishes the proposed matching.  These details, rather than a general
strong-unitarity implication, are the content of Claim~A1.

\begin{theorem}[Conditional complete unitarity]
\label{thm:analytic-complete-unitarity}
Assume Claims~\ref{clm:analytic-net-count},
\ref{clm:analytic-reconstruction}, and~\ref{clm:analytic-annulus},
together with the preceding algebraic inputs.  Then $U$ is completely
unitary.
\end{theorem}
\begin{proof}
Proposition~\ref{prop:analytic-strong-unitarity} unitarizes every simple
module.  Claim~\ref{clm:analytic-annulus} supplies the representation
required by
\href{https://link.springer.com/article/10.1007/s00039-024-00685-8}
{Tener, Definition~2.18}, for each such module; finite direct sums
preserve the property.  His Corollary~4.6 then supplies positivity of
the canonical fusion forms.  Strong unitarity must be established first:
the unitarizable-fusion-product hypothesis in Theorem~4.4 cannot be used
to bypass the exhaustion argument.
\end{proof}

\section{Conditional unitary categorical identification}
\label{sec:analytic-category}

\begin{theorem}[Unitary Haagerup-center identification, conditional on A1--A3]
\label{thm:analytic-category}
Under the hypotheses of Theorem~\ref{thm:analytic-complete-unitarity},
and the preceding determination of the all-module fusion ring and its
untwisted pointed subcategory, there is a unitary ribbon equivalence
\[
 \Rep(U)\simeq Z(\HH_3).
\]
For the original embedded extension $H/U$, its all-module category is
tensor equivalent to one of $\HH_2$ and $\HH_3$.
\end{theorem}
\begin{proof}
Put $\mathcal C=\Rep(U)$ and let $A$ be the algebra of the actual
unitary extension $H$.  Complete unitarity allows application of
\href{https://arxiv.org/html/1908.03484#S2.SS4}
{Gui, Theorems~2.21 and~2.28}: $A$ is a standard commutative
Q-system and its modules are unitarizable.  The unitary braiding gives
each one-sided module a compatible opposite-side action, realizing
$\mathcal F=\mathcal C_A$ in the unitary bimodule construction.  Thus
$\mathcal F$ is a $C^*$ fusion category.  Here $\mathcal C_A$ means
\emph{all} modules.  The local modules form $\mathcal C_A^0\simeq\Vec$
because the extension is holomorphic, so substituting local modules
would lose the Haagerup category.

Its already determined fusion ring is
$\operatorname{HI}(\mathbb Z/3)$.  The actual untwisted pointed
subcategory makes the obstruction $c^{(0,3)}$ vanish, while the other
obstruction is in
\[
 H^1\!\left(\mathbb Z/2,H^2(\mathbb Z/3,\mathbb C^\times)\right)=0.
\]
These are the hypotheses in
\href{https://arxiv.org/html/1609.07604}
{Izumi, Definition~2.7 and Theorems~5.9 and~9.1}, giving the two
unitary tensor classes $\HH_2,\HH_3$.  No algebra structure on
$1+\rho$ is inserted as an extra assumption.
\href{https://arxiv.org/html/1102.2631#S1}
{Grossman--Snyder, Theorem~1.4}, places these categories in the same
Morita class.  Their centers are therefore braided equivalent.
The positive spherical structures make this a ribbon equivalence, and
the Lagrangian-center theorem gives
\[
 \mathcal C\simeq Z(\mathcal F)\simeq Z(\HH_3).
\]
\href{https://arxiv.org/html/1906.09710}
{Reutter, Theorem~5}, provides the unitary braided realization of the
equivalence, preserving the associated positive ribbon structures.
\end{proof}

This uses the published unitary classification as an external input;
its classification equations have not been replayed here.  It avoids
the broader pivotal census, but it does not prove that the natural
reconstruction functor $\mathcal R$ preserves Connes fusion, the VOA
tensor product, or braiding.  That particular comparison is a distinct
statement and is not needed for the theorem just proved conditionally.

\section{Conditional full theories and a chosen \texorpdfstring{$E_8$}{E8} extension}
\label{sec:analytic-full-theory}

\begin{proposition}[Diagonal full theory; conditional character identification]
\label{prop:analytic-diagonal}
The completely rational net $\mathcal N_U$ admits its canonical unitary
local diagonal full theory.  Under Claims~\ref{clm:analytic-net-count}
and~\ref{clm:analytic-reconstruction}, and the preceding twelve-module
character determination, its partition function is
\begin{equation}
 Z_{\mathrm{diag}}(\tau,\bar\tau)
   =\sum_{i=0}^{11}\chi_i(\tau)\overline{\chi_i(\tau)},
 \qquad c_L=c_R=8.
 \label{eq:analytic-diagonal-partition}
\end{equation}
\end{proposition}
\begin{proof}
Apply the diagonal Longo--Rehren construction, in the identity-coupling
case of
\href{https://arxiv.org/html/math/9911149}
{Rehren, Theorem~1.4, Proposition~1.5, and Corollary~1.6}.
Claims A2--A3 identify all sector Hilbert spaces with the twelve ordinary
modules and preserve their Hamiltonians.  Hence their traces give
\eqref{eq:analytic-diagonal-partition}.  The annulus compatibility claim
is not needed for this net construction or this linear character matching.
\end{proof}

\begin{theorem}[Chosen $\HH_3$ extension and defects; conditional on A1--A3]
\label{thm:analytic-e8-extension}
Assume the hypotheses of Theorem~\ref{thm:analytic-category} and the
preceding strict positivity of all nonvacuum lowest conformal weights.
There exists a unitary holomorphic conformal extension $H'\supset U$
isomorphic to $V_{E_8}$ whose all-module category is $\HH_3$.
For its extension algebra $A_3\in\Rep(U)$,
\begin{equation}
 \Rep(U)_{A_3}\simeq\HH_3,
 \qquad
 {}_{A_3}\Rep(U)_{A_3}\simeq\HH_3\boxtimes\HH_3^{\mathrm{op}}.
 \label{eq:analytic-extension-defects}
\end{equation}
The second category describes algebra-bimodule defects preserving the
$U$ chiral subalgebra in the corresponding full extension theory and
contains an $\HH_3$ fusion subcategory.
\end{theorem}
\begin{proof}
Let $I_{\HH_3}$ be the right adjoint of the forgetful functor
$Z(\HH_3)\to\HH_3$.  Transport the canonical algebra
$I_{\HH_3}(\mathbf1)$ along the unitary ribbon equivalence.  By the
canonical Lagrangian-algebra construction, as in
\href{https://arxiv.org/abs/1009.2117}
{Davydov--M\"uger--Nikshych--Ostrik, Section~4}, the resulting $A_3$
is connected, commutative, separable, and twist-free, with
\[
 \FPdim(A_3)^2=\FPdim(\Rep(U)),
 \qquad \Rep(U)_{A_3}^{0}\simeq\Vec,
 \qquad \Rep(U)_{A_3}\simeq\HH_3.
\]
Trivial twist makes its constituent conformal weights integral.
Connectedness and strict positivity away from the vacuum give a
one-dimensional weight-zero space.  The extension correspondence of
\href{https://arxiv.org/html/1406.3420}
{Huang--Kirillov--Lepowsky, Theorems~3.2--3.6}, together with Gui's
unitary Q-system theorem, produces a unitary conformal VOA extension
$H'$.  It is a finite sum of ordinary $U$-modules; finite-extension
$C_2$-cofiniteness and the semisimple local-module description give
strong rationality.  The Lagrangian condition makes it holomorphic.
Its central charge is eight, so
\href{https://msp.org/pjm/2004/213-2/pjm-v213-n2-p05-p.pdf}
{Dong--Mason, Theorem~1}, identifies $H'$ with $V_{E_8}$.

For the canonical algebra, the algebra-bimodule duality gives the
second equivalence in \eqref{eq:analytic-extension-defects}; see
\href{https://arxiv.org/abs/math/0111205}
{M\"uger, Sections~4 and~8}.  Here ``op'' reverses the tensor product.
The first-factor embedding $x\mapsto x\boxtimes\mathbf1$ supplies
the stated $\HH_3$ fusion subcategory.
\end{proof}

\begin{remark}[Extension choice and preserved symmetry]
\label{rem:analytic-extension-scope}
The selection of the canonical $\HH_3$ algebra proves existence of
$H'$ with the specified all-module category.  It does not identify
$A_3$ with the original embedded $H/U$ algebra; the original boundary
category may be $\HH_2$ or $\HH_3$.  Both holomorphic VOAs being
isomorphic to $V_{E_8}$ does not identify their embeddings of $U$.
Moreover, the defects in \eqref{eq:analytic-extension-defects} preserve
$U$ and need not preserve the enlarged $E_8$ chiral algebra.
The local chiral $H'$-module category is $\Vec$.
\end{remark}

The logical separation is therefore explicit: completed finite
certificates prove the actual strong-generation and properness result;
the algebraic identification supplies its stated representation-theory
inputs; A2--A3 supply sector exhaustion and strong unitarity; A1 supplies
the additional analytic comparison needed for complete unitarity.
The unitary categorical and $\HH_3$ extension conclusions retain those
conditions until the highlighted applications are independently verified.
The synthesis sources are
\artifact{hi_cft_20260923_c8two_final_theorem.md} and
\artifact{hi_cft_20260923_c8two_certificate_audit.md}; their unqualified
analytic conclusion is deliberately recorded here with the review boundary
made visible.


\section{Comparison with arXiv:2512.08225}\label{sec:comparison}
Albert--Honda--Kaidi--Zheng explicitly assume the existence of a
\(c=8\) Haagerup-center theory in Section II, footnote 4 of
\href{https://arxiv.org/html/2512.08225v2#S2}{arXiv:2512.08225v2}.
Their equation (12) displays the \(\gamma\)-dependent characters,
and their extension discussion relates the center theory to the
\(E_8\) theory with Haagerup defects.

If the highlighted claims of this manuscript are confirmed, the explicit
algebra \eqref{eq:U} supplies their \(\gamma=0\) existence premise.
The choice of an \(\HH_3\) Lagrangian algebra supplies a compatible
\(E_8\) extension, with the distinction between \(U\)-preserving and
full-\(E_8\)-preserving defects maintained. Neither this construction
nor that finite gauging identifies the continuum limit of the
\(\rho\)-projector chain. Section III of the cited paper proposes
further asymmetric quotients towards charges two and six; those require
additional conformal and defect data. No uniqueness theorem for all
central-charge-eight CFTs is asserted here.

\appendix
\section{Certificate ledger and reproducibility}\label{sec:ledger}
The following are existing, closed records. No new lattice computation
was run while preparing this typeset manuscript. All paths are relative
to the project directory. The long filenames deliberately preserve an
unambiguous connection to the actual records.

\begin{longtable}{@{}p{0.25\linewidth}p{0.67\linewidth}@{}}
\toprule
Result & Record and scope\\
\midrule\endfirsthead
\toprule Result & Record and scope\\\midrule\endhead
\bottomrule\endfoot
Lower identities &
\artifact{hi_cft_20260923_next_structure_lower_closure.json}:
3,970 exact relations.\\
Actual triplet &
\artifact{hi_cft_20260923_next_structure_p4_controls.json}:
18 current-zero and six direct Virasoro controls.\\
Raising saturation &
\artifact{hi_cft_20260923_c8two_structure_saturation.json}:
644 exact images in a quotient of dimension 322.\\
Additional direct seeds &
\artifact{hi_cft_20260923_c8two_targeted_h4_trimmed_v2/summary.json}:
118 completed exact identities. The unfinished final block in this
historical directory is not counted as a completed identity.\\
Raised seeds &
\artifact{hi_cft_20260923_c8two_structure_h4_raised_exact/summary.json}:
twelve additional ordinary exact identities.\\
Singlet premise &
\artifact{hi_cft_20260923_c8two_h4_singlet_premise_verified.json}:
all 24 raising conditions replayed.\\
Projected residuals &
\artifact{hi_cft_20260923_c8two_h4_singlet_exact.json}:
twelve full exact residuals, using 21 selected columns.\\
Full Gate 4 &
\artifact{hi_cft_20260923_c8two_gate4_verified.json}:
all 2,898 identities, using singlet recovery and raising saturation.\\
Charged spanning &
\artifact{hi_cft_20260923_c8two_charged_span_verified.json}:
28 reconstructed seeds, rank 322, and current commutators.\\
Full Gate 5 &
\artifact{hi_cft_20260923_c8two_gate5_verified.json}:
all 24 channels by the stated Jacobi argument.\\
Combined provenance &
\artifact{hi_cft_20260923_c8two_proof_bundle.json}:
640 authenticated files, including earlier exact inputs and proof notes.
It checks provenance and completion, not formal correctness of every
mathematical argument.\\
\end{longtable}

The final Gate 4 verifier replays the twelve projected identities using
integer-cleared coefficients, separately from the rational producer.
A field element is represented in the basis
\(1,s,t,st\), with \(s^2=13\), \(t^2=-3\); all residuals are
exact four-component rational identities. Prime calculations are used
to select candidate columns and to certify positive spanning against
an independently known upper dimension. A vanishing prime residual
alone is never treated as a characteristic-zero containment proof.

The provenance-ledger replay can be written to a fresh output without
changing the original certificate:
\begin{quote}\small
\begin{verbatim}
python hi_cft_safe_run.py --limit-mib 1024 -- \
  python hi_cft_20260923_c8two_proof_bundle.py \
  --output /tmp/haagerup-c8-proof-review.json
\end{verbatim}
\end{quote}
These lines represent a single shell command. The output path must not
already exist. This rechecks the frozen source hashes and completion
flags; it does not rerun all original mode calculations. The source
files of the individual certificate producers and verifiers are bound
into the same ledger.

\section{Compilation and review scope}\label{sec:compilation}
The main source is \artifact{hi_cft_c8_proof.tex}; its companion sources
are \artifact{hi_cft_c8_manuscript_structure.tex} and
\artifact{hi_cft_c8_manuscript_analytic.tex}. Run \texttt{pdflatex} twice
on the main source for cross-references and the table of contents.
The document uses standard TeX packages and requires no shell escape.

The computational records predate this manuscript and remain unchanged.
The original provenance ledger consequently does not authenticate this
new exposition. The PDF organizes and qualifies those records; it is not
itself an additional mathematical certificate. Before treating the full
identification as externally established, the highlighted local claims
should be checked by specialists and the certificate recipes should be
replayed in an independently reviewed environment.
\end{document}
